% chapters/conclusions/achieved_objectives.tex
% Sección 7.1: Conclusiones (~1-1.5 páginas)
% ============================================================================
%
% GUÍA: Repasar los objetivos del \cref{sec:objectives} uno a uno y valorar
%       su grado de consecución. Es MUY importante que cada objetivo (O-1,
%       O-2, etc.) se mencione explícitamente aquí.
%
%   Estructura recomendada:
%
%   1. Párrafo introductorio: "En el \cref{sec:objectives} se plantearon
%      N objetivos para este TFG. A continuación se valora el grado de
%      consecución de cada uno."
%
%   2. Para cada objetivo:
%      "El \cref{obj:api} (O-1) planteaba diseñar e implementar una API REST...
%       Este objetivo se ha alcanzado satisfactoriamente, como demuestran
%       los N endpoints implementados y las pruebas de integración del
%       \cref{sec:integration_tests}."
%
%      Usar \ref{obj:xxx} para referenciar cada objetivo.
%
%   3. Párrafo de cierre: valoración global del trabajo, dificultades
%      principales, lecciones aprendidas.
%      - Qué fue lo más difícil (p.ej., pipeline OCR, máquina de estados).
%      - Qué se aprendió (tecnologías, ingeniería de software, gestión).
%      - Satisfacción con el resultado.
%
%   NOTA: No introducir información nueva aquí; solo referencias a
%   capítulos anteriores.
% ============================================================================

% TODO: Redactar las conclusiones (se hace al final, cuando todo esté hecho).
